\documentclass[10pt,twocolumn]{article}
\usepackage[T1]{fontenc}
\usepackage[utf8]{inputenc}
\usepackage{lmodern}
\usepackage[margin=1.8cm,top=2.4cm,bottom=2cm,columnsep=0.6cm]{geometry}
\usepackage{fancyhdr}
\usepackage{titlesec}
\usepackage{booktabs}
\usepackage{amsmath,amssymb,amsfonts}
\usepackage{microtype}
\usepackage{xcolor}
\usepackage{mdframed}
\usepackage{parskip}
\usepackage{enumitem}
\usepackage{hyperref}

\definecolor{rulered}{RGB}{180,30,30}
\definecolor{boxgray}{RGB}{245,245,245}
\definecolor{keywordgray}{RGB}{100,100,100}

\pagestyle{fancy}
\fancyhf{}
\fancyhead[L]{\textsc{\small Claude Mag}}
\fancyhead[C]{\small\textit{Machine Learning Research Digest}}
\fancyhead[R]{\small 2026-09-14}
\fancyfoot[C]{\small\thepage}
\renewcommand{\headrulewidth}{0.8pt}

\titleformat{\section}{\normalsize\bfseries\color{rulered}}{}{0pt}{}%
  [\vspace{-4pt}\color{rulered}\rule{\columnwidth}{0.4pt}]
\titlespacing{\section}{0pt}{8pt}{4pt}

\hypersetup{colorlinks=true,urlcolor=rulered,linkcolor=black}

\begin{document}
\setlength{\parindent}{0pt}
\setlength{\parskip}{4pt}

{\small\color{keywordgray}\textbf{Keywords:}
  causal discovery \textbullet{}
  foundation models \textbullet{}
  benchmark \textbullet{}
  structural causal models \textbullet{}
  memorization}\par\vspace{4pt}

{\LARGE\bfseries\raggedright
  CausalArena: Benchmarking Causal
  Discovery in the Foundation Model Era}\par\vspace{4pt}

{\large\itshape\raggedright
  Fixed Synthetic Benchmarks Overestimate CDFM Performance by
  15--28 Points Due to Pretraining Overlap ---
  Classical Algorithms Remain Competitive on Novel SCMs}\par\vspace{6pt}

{\small\textbf{By} Zi-Rong Li, Si-Yang Liu,
  Tian-Zuo Wang, Han-Jia Ye
  (Nanjing University)}\par\vspace{2pt}

{\small\color{keywordgray}arXiv:2609.11897}
\par\vspace{2pt}\noindent{\color{rulered}\rule{\columnwidth}{0.6pt}}\vspace{6pt}

Causal discovery has entered the foundation model era: pretrained
causal discovery foundation models (CDFMs) learn to predict DAGs
in-context, and LLM-based causal agents reason about structure from
variable names. Both paradigms achieve impressive scores on existing
benchmarks --- but these benchmarks are fixed, synthetic, and drawn
from the same SCM families that CDFMs were pretrained on. CausalArena
introduces a unified, evolvable benchmark that controls for pretraining
overlap and enables, for the first time, meaningful cross-paradigm
comparisons. On genuinely novel SCMs, classical algorithms match or
outperform CDFMs; on semantically rich named-variable graphs, LLM
agents dominate. No paradigm is universally best.

\begin{mdframed}[backgroundcolor=boxgray,linecolor=rulered,linewidth=1pt,
    innerleftmargin=6pt,innerrightmargin=6pt,
    innertopmargin=6pt,innerbottommargin=6pt]
\textbf{\small AT A GLANCE}\par\vspace{4pt}
\begin{description}[leftmargin=0pt,labelsep=4pt,itemsep=2pt,parsep=0pt,
    font=\small\bfseries]
  \item[Problem:] Fixed synthetic benchmarks inflate CDFM performance
    due to pretraining--test overlap; cross-paradigm comparison is
    impossible under inconsistent protocols.
  \item[Why it matters:] Practitioners need to choose the right
    causal discovery tool; misleading benchmarks produce wrong choices.
  \item[Method:] Procedural SCM generator producing novel graphs;
    common protocol with SHD, F1, and new CER metric; evolvable
    test set refreshed monthly.
  \item[Result:] CDFMs score 8--14 SHD points worse on out-of-pretraining
    SCMs; classical algorithms match CDFMs on novel SCMs; LLM agents
    dominate on named-variable graphs.
  \item[Evidence:] 11 algorithms across 3 paradigms; 12 SCM families;
    canary memorization test; 3 data regimes.
\end{description}
\end{mdframed}

\section{The Evaluation Crisis in Causal Discovery}

Causal discovery is the task of inferring a directed acyclic graph
(DAG) $\mathcal{G} = (V,E)$ from observational data
$\mathbf{X} \sim P_\mathcal{G}$. Three paradigms compete:

\textbf{Classical algorithms} (PC, FCI, GES, NOTEARS) operate
directly on data via conditional independence tests or structural
score optimization. They have no pretraining.

\textbf{CDFMs} (CausalFormer, DiffAN, CDFL) are pretrained on
millions of synthetic SCM-generated datasets and predict DAGs
in-context given a new dataset --- analogous to in-context learning
in language models.

\textbf{LLM causal agents} (GPT-4o-Causal, Claude-3.5-Causal) use
semantic knowledge of variable names and domain background to reason
about causal relations, augmented by statistical tests.

Each paradigm has been evaluated on different benchmarks with different
graph families, mechanisms, and metrics. CDFMs report large wins over
classical algorithms on their benchmark --- but their benchmark uses
the same SCM families they were pretrained on.

\section{The Memorization Problem}

The authors prove that CDFMs exhibit significant memorization: their
performance on SCM families that appear in their pretraining corpus
is 15--28 SHD points better than on structurally novel families.

\textbf{Memorization gap:}
\[
  \Delta_\mathrm{mem} = \mathrm{SHD}(M,\, \mathcal{G}_\mathrm{canonical})
    - \mathrm{SHD}(M,\, \mathcal{G}_\mathrm{canary})
\]
where $\mathcal{G}_\mathrm{canary}$ is a ``canary'' SCM family
constructed by the authors that does not appear in any published
CDFM pretraining curriculum. All three CDFMs show
$\Delta_\mathrm{mem} > 5$ SHD points ($p < 0.001$).

This means published CDFM benchmark numbers are inflated by
pretraining memorization, not causal discovery ability.

\section{CausalArena Design}

\textbf{SCM Generator.} CausalArena procedurally generates novel
SCMs outside any existing CDFM pretraining corpus:
\begin{itemize}[itemsep=2pt,topsep=2pt]
  \item \textit{Graph families:} Erd\H{o}s--R\'enyi, scale-free,
    bipartite layered, causal motif (new).
  \item \textit{Mechanisms:} Linear Gaussian, ANM, post-nonlinear,
    discretized multinomial.
  \item \textit{Semantics:} Anonymized vs.\ domain-named variables.
  \item \textit{Sizes:} 5--50 nodes.
\end{itemize}

\textbf{Common Protocol.} CausalArena enforces:
\begin{itemize}[itemsep=2pt,topsep=2pt]
  \item \textit{Metrics:} SHD, F1 on directed edges, and the new
    Causal Effect Recovery (CER) metric.
  \item \textit{Data regimes:} $n \in \{500, 5000, 50000\}$.
  \item \textit{Knowledge levels:} Zero-knowledge (anonymous),
    partial (names), full (names + background document).
\end{itemize}

\textbf{Evolvability.} The test set is seeded by a public random key
and refreshed monthly to prevent re-memorization. A version-controlled
leaderboard tracks progress over time.

\section{The CER Metric}

SHD measures graph structure but not causal effect accuracy. CausalArena
introduces \textbf{Causal Effect Recovery (CER)}:
\[
  \mathrm{CER} = 1 - \frac{1}{|P|} \sum_{(i,j) \in P}
    \frac{|\hat{\tau}_{ij} - \tau_{ij}|}{|\tau_{ij}|}
\]
where $P$ is the set of all variable pairs, $\tau_{ij}$ is the true
total causal effect of $X_i$ on $X_j$ (via $do$-calculus), and
$\hat{\tau}_{ij}$ is estimated from the predicted DAG via structural
equation fitting. CER directly measures whether the predicted graph
supports correct downstream interventional reasoning.

\section{Experimental Findings}

\textbf{Novel SCMs, $n=5{,}000$, anonymized variables:}

\begin{center}
\begin{tabular}{lccc}
\toprule
Algorithm & SHD $\downarrow$ & F1 $\uparrow$ & CER $\uparrow$ \\
\midrule
GES              & 7.9  & 0.64 & 0.74 \\
PC               & 8.4  & 0.61 & 0.71 \\
CausalFormer     & 9.8  & 0.57 & 0.65 \\
DiffAN           & 10.3 & 0.54 & 0.63 \\
GPT-4o-Causal    & 18.7 & 0.31 & 0.41 \\
\bottomrule
\end{tabular}
\end{center}

Classical algorithms outperform CDFMs on novel, anonymized SCMs ---
reversing the apparent superiority CDFMs show on fixed benchmarks.

\textbf{Named variables, $n=5{,}000$:}

\begin{center}
\begin{tabular}{lccc}
\toprule
Algorithm & SHD $\downarrow$ & F1 $\uparrow$ & CER $\uparrow$ \\
\midrule
GES                  & 8.1 & 0.63 & 0.73 \\
CausalFormer         & 10.1 & 0.56 & 0.64 \\
\textbf{GPT-4o-Causal} & \textbf{6.2} & \textbf{0.74} & \textbf{0.82} \\
\bottomrule
\end{tabular}
\end{center}

LLM agents dominate when variable names carry domain-specific semantics.

\section{Ablations and Interpretation}

\textbf{Sample size:} CDFMs show flat performance across
$n \in \{500, 5000, 50000\}$ --- they rely on learned structure priors.
Classical algorithms improve monotonically with $n$.

\textbf{Mechanism complexity:} Post-nonlinear mechanisms are hardest
for all methods; CDFMs trained on linear Gaussian data show a
12-point SHD gap on nonlinear instances.

\textbf{Graph size:} All algorithms degrade with more nodes; classical
algorithms degrade more gracefully at $n > 20$ nodes than CDFMs.

\textbf{Semantic grounding:} LLM agents degrade sharply on anonymized
variables ($-12.5$ SHD points), showing their advantage is entirely
domain-knowledge driven.

\section*{Reference}

Zi-Rong Li, Si-Yang Liu, Tian-Zuo Wang, Han-Jia Ye.
\textbf{CausalArena: Benchmarking Causal Discovery in the Foundation
Model Era.} arXiv:2609.11897, September 2026.
\url{https://arxiv.org/abs/2609.11897}

\end{document}
